\begin{env}[7.9.2]
\label{IV.2.s7.9.2-fr}
Es claro que, si se puede resolver $X$, también se puede resolver todo
preesquema inducido en un \emph{abierto} de $X$ y todo preesquema local
$\operatorname{Spec}(\mathcal O_x)$ de $X$ (\textbf{II}, 5.4.2).
Por otra parte, si se puede resolver $X$, es claro que existe en $X$
un abierto \emph{denso en todo el preesquema} contenido en $\operatorname{Reg}(X)$.
Estas observaciones muestran inmediatamente que, si para todo
subpreesquema \emph{cerrado e íntegro} $Y$ de $X$ y todo $Y$-preesquema
$Y'$ \emph{íntegro, finito} y \emph{radicial} sobre $Y$ se puede resolver
$Y'$, entonces los abiertos afines de $X$ satisfacen \emph{la condición}
(iii) \emph{de} (7.8.2).

Por otra parte, para los esquemas locales se tiene la proposición siguiente:
\end{env}
